상업적 민감성 때문에, 저자들은 다운스트림 절대 지표를 보고하지 않고 모든 결과를 과제별 기준 대비 상대 변화로 표기한다.
본 논문 전반에서 상대 성능은 $(x / \text{baseline} - 1)\,\%$로 계산되며, 여기서 $x$는 평가 대상 방법의 점수다.

\subsection{평가 프로토콜}
\label{sec:protocol}

저자들은 PRAGMA를 주로 다운스트림 과제에 대한 임베딩 프로브와 Low-Rank Adaptation(LoRA)~\citep{hu2022lora} 파인튜닝으로 평가한다.

\subsubsection{임베딩 프로빙}
임베딩 프로빙은 LoRA 파인튜닝까지 가지 않고도 실험 반복을 빠르게 돌릴 수 있게 해 준다 — 예컨대 새 피처가 기대만큼 이득을 가져다주는지 가늠하거나, 사전학습 실행 후 추가 평가를 위한 체크포인트를 고르거나, 어떤 과제를 다운스트림 타깃으로 탐색할 가치가 있는지 결정하는 데 쓰인다.
임베딩은 히스토리 인코더 출력($z_h$)에서 추출한다.

프로빙 분석에서 저자들은 \texttt{[USR]} 토큰, 마지막 \texttt{[EVT]} 토큰, 그리고 둘의 결합을 표준 선형 프로브로 평가한다.
미리 정의된 학습/검증/테스트 분할이 있는 다운스트림 과제에서, 먼저 각 레코드를 동결된 인코더로 통과시켜 고정 크기 표현을 얻고, 학습 분할에서 선형 프로브(로지스틱 또는 선형 회귀)를 학습한다.
프로브 성능은 하이퍼파라미터 선택에 강건하기 때문에, 프로브 적합은 보통 수 분이면 끝난다.
저자들의 아키텍처가 본질적으로 ``pre-norm''이므로, 프로브 적합 전에 임베딩에 표준 스케일링을 적용한다.
프로브 학습에는 L-BFGS 옵티마이저~\citep{liu1989limited}가 가장 좋은 결과를 내고 빠르게 수렴함을 확인했다.

참고로, 차원이 낮은 임베딩(예: $192$차원)에서는 그래디언트 부스팅 결정 트리(GBDT)도 잘 작동하지만, 과제별 하이퍼파라미터 튜닝이 필요하고 적합 시간이 더 길어, 빠른 모델 평가용으로는 선형 프로빙보다 덜 실용적이다.

\subsubsection{LoRA를 이용한 다운스트림 적응}
PRAGMA 백본을 다운스트림 과제에 특화시키기 위해, 저자들은 단지 2--4\,\%의 파라미터 오버헤드만 추가하는 Low-Rank Adaptation(LoRA)을 사용한다.
이 설정에서는 사전학습된 가중치를 과제별 목적함수에 맞춰 파인튜닝함으로써, 일반 표현 학습과 다운스트림 요구 사이의 간극을 메운다.

LoRA는 일반 관행을 따라 인코더 레이어 안의 QKV 사영과 MLP 레이어에 적용한다~\citep{hu2022lora,dettmers2023qlora}. 모든 실험에서 기본값으로 $\text{rank}=8$, $\alpha=8$을 쓰지만, 더 작은 데이터셋에서는 $\{4, 8, 16\}$ 범위로 rank를 스윕한다.
LoRA 파인튜닝에는 Adam 옵티마이저~\citep{kingma2015adam}를 사용하며, 학습은 보통 사전학습 벽시계 시간의 약 1/8을 쓰고, 데이터셋 크기에 따라 12~시간에서 며칠 사이에 수렴한다.

\subsubsection{다운스트림 데이터셋 준비}
각 다운스트림 과제에 대해 저자들은 고유 식별자를 얻는다 — 보통 프로필 id와 평가 시점으로 구성된다.
다음으로, 평가 시점 직전까지의 이벤트 이력과 프로필 속성을 모은다.
각 다운스트림 과제에 대해 미리 정해진 폴드와 분할을 따른다.
다운스트림 데이터셋 수집 절차는 사전학습 데이터셋의 절차를 그대로 반영한다.
